\section{Longest Flight Route}
\label{sec:cses-longest-flight-route}

\problemstatement{Given a \emph{directed acyclic} graph of $n$ cities and
$m$ flights, find a route from city $1$ to city $n$ visiting the maximum
number of cities, and print it, or \texttt{IMPOSSIBLE}.}

\begin{patternbox}
Longest path in a \textbf{DAG} is a topological-order DP -- unlike general
graphs (where longest path is NP-hard), acyclicity makes it linear. Process
nodes so that every predecessor is finalised first, and each node takes the
best incoming route.
\end{patternbox}

\subsection*{DP over topological order}

Let $\textit{dp}[v]$ be the maximum number of cities on a route from $1$ to
$v$. Initialise $\textit{dp}[1]=1$ and others $-\infty$ (unreachable).
Processing nodes in topological order, relax each edge $u\to v$: if
$\textit{dp}[u]+1>\textit{dp}[v]$, set $\textit{dp}[v]=\textit{dp}[u]+1$ and
record $u$ as $v$'s predecessor. Because $u$ precedes $v$ topologically,
$\textit{dp}[u]$ is final when used. Reconstruct the route by following
predecessors back from $n$.

\begin{ideabox}[Acyclicity Turns Longest Path into DP]
In a DAG the topological order guarantees each node's best value is computed
before it is needed downstream, so a single sweep finds the longest path.
The predecessor array turns the length into a printable route.
\end{ideabox}

\subsection*{Algorithm}

\begin{enumerate}
  \item Topologically sort (Kahn's). $\textit{dp}[1]=1$, rest $-\infty$.
  \item In topological order, relax out-edges, updating $\textit{dp}$ and
        predecessors.
  \item If $\textit{dp}[n]<0$ print \texttt{IMPOSSIBLE}; else backtrack
        from $n$ and print the route.
\end{enumerate}

\subsection*{Dry run}

\dryrun{$1\!\to\!2$, $1\!\to\!3$, $2\!\to\!3$, $3\!\to\!4$; $n=4$.}

\begin{itemize}
  \item $\textit{dp}$: $1{:}1,2{:}2,3{:}3$ (via $1\!\to\!2\!\to\!3$),
        $4{:}4$.
  \item Longest route $1\to2\to3\to4$, four cities.
\end{itemize}

\subsection*{C++ implementation}

\begin{lstlisting}
#include <bits/stdc++.h>
using namespace std;

int main() {
    int n, m;
    cin >> n >> m;
    vector<vector<int>> adj(n + 1);
    vector<int> indeg(n + 1, 0);
    for (int i = 0; i < m; i++) {
        int a, b; cin >> a >> b;
        adj[a].push_back(b);
        indeg[b]++;
    }

    // topological order via Kahn's
    vector<int> order, deg = indeg;
    queue<int> q;
    for (int i = 1; i <= n; i++) if (deg[i] == 0) q.push(i);
    while (!q.empty()) {
        int u = q.front(); q.pop(); order.push_back(u);
        for (int v : adj[u]) if (--deg[v] == 0) q.push(v);
    }

    const int NEG = -1e9;
    vector<int> dp(n + 1, NEG), pre(n + 1, 0);
    dp[1] = 1;
    for (int u : order)
        if (dp[u] > NEG)
            for (int v : adj[u])
                if (dp[u] + 1 > dp[v]) { dp[v] = dp[u] + 1; pre[v] = u; }

    if (dp[n] < 0) { cout << "IMPOSSIBLE\n"; return 0; }
    vector<int> route;
    for (int v = n; v != 0; v = pre[v]) route.push_back(v);
    reverse(route.begin(), route.end());
    cout << route.size() << "\n";
    for (int v : route) cout << v << " ";
    cout << "\n";
    return 0;
}
\end{lstlisting}

\begin{complexitybox}
\textbf{Time Complexity.} $O(n+m)$ -- topological sort plus one relaxation
sweep. \textbf{Space Complexity.} $O(n+m)$.
\end{complexitybox}

\begin{takeawaybox}
Longest path is intractable in general but linear on a DAG: relax edges in
topological order and track predecessors. The acyclicity guarantee is the
entire reason this DP is valid -- the same guarantee Game Routes (next)
exploits for path counting.
\end{takeawaybox}
